% GENERATED FILE. Edit canonical lessons, then run npm run generate:books.
% canonical-source-hash: fnv1a64:fcfb7920382bac90
% canonical-lessons: ML-C438-grahikkuka, ML-C438-hrdisthamakkuka, ML-C438-parisilippikkuka, ML-C438-vyakhyanikkuka, ML-C438-nirvacikkuka

\chapter{Ways of Learning}
\label{ch:ml-ways-of-learning}
\hlchaptermodality{438}

\begin{chapteropening}
\textbf{By the end of this chapter:} \emph{I can say to grasp, to memorise, to train, to interpret and to define.}

Complete the last lesson of chapter 438: I can say to grasp, to memorise, to train, to interpret and to define.
\end{chapteropening}

\section[\texorpdfstring{\ml{ഗ്രഹിക്കുക} (grahikkuka)}{ഗ്രഹിക്കുക (grahikkuka)}]{\ml{ഗ്രഹിക്കുക} (grahikkuka) — to grasp; to comprehend}

\label{lesson:ML-C438-grahikkuka}



\begin{quote}
Before the new one: say the Malayalam for knowledge, then the Malayalam for studies.
\end{quote}

\subsection*{You'll want to know: \ml{ഗ്രഹിക്കുക}}
\textbf{\ml{ഗ്രഹിക്കുക}} — \emph{grahikkuka} — \textquotedblleft{}to grasp; to comprehend\textquotedblright{}.

\begin{quote}
\ml{അവൾ} \ml{ആ} \ml{പാഠം} \ml{വേഗം} \ml{ഗ്രഹിച്ചു}.
\end{quote}

\begin{quote}
\emph{avaḷ ā pāṭhaṁ vēgaṁ grahiccu.}
\end{quote}

\begin{quote}
She grasped that lesson quickly.
\end{quote}

Past: \ml{ഗ്രഹിച്ചു}, \emph{grahiccu}.

A bookish word; in talk people say \ml{മനസ്സിലാക്കുക}.

\begin{cousinweb}
\textbf{\ml{ഗ്രഹിക്കുക}} (\emph{grahikkuka}) is built from Sanskrit \emph{grahaṇa}, \textquotedblleft{}seizing, holding\textquotedblright{} (CDIAL 4364), with the Malayalam verb ending -\ml{ിക്കുക}. To grasp an idea is to seize it, in Malayalam as in English.
\end{cousinweb}

\begin{grammarlens}[title={What you've built}]
One more word for giving reasons and opinions.
\end{grammarlens}

\subsection*{Your turn}
\begin{itemize}
\raggedright
  \item \emph{Say it:} \emph{grahikkuka}
  \item \emph{Say it:} \emph{grahikkuka}, once more
  \item \emph{Say it:} \emph{avaḷ ā pāṭhaṁ vēgaṁ grahiccu.} — She grasped that lesson quickly.
  \item \emph{Recall:} say the Malayalam for knowledge, then the Malayalam for studies, then say \emph{grahikkuka} again
\end{itemize}

\begin{tcolorbox}[breakable,colback=teal!4,colframe=teal!35!black,title={Before you move on}]
What does \ml{ഗ്രഹിക്കുക} mean? (\textquotedblleft{}To grasp; to comprehend\textquotedblright{}.) Say it once more.
\end{tcolorbox}
\section[\texorpdfstring{\ml{ഹൃദിസ്ഥമാക്കുക} (hṛdisthamākkuka)}{ഹൃദിസ്ഥമാക്കുക (hṛdisthamākkuka)}]{\ml{ഹൃദിസ്ഥമാക്കുക} (hṛdisthamākkuka) — to memorise; to learn by heart}

\label{lesson:ML-C438-hrdisthamakkuka}



\begin{quote}
Before the new one: say the Malayalam for studies, then the Malayalam for to grasp.
\end{quote}

\subsection*{You'll want to know: \ml{ഹൃദിസ്ഥമാക്കുക}}
\textbf{\ml{ഹൃദിസ്ഥമാക്കുക}} — \emph{hṛdisthamākkuka} — \textquotedblleft{}to memorise; to learn by heart\textquotedblright{}.

\begin{quote}
\ml{കുട്ടി} \ml{ആ} \ml{കവിത} \ml{മുഴുവൻ} \ml{ഹൃദിസ്ഥമാക്കി}.
\end{quote}

\begin{quote}
\emph{kuṭṭi ā kavita muḻuvan hṛdisthamākki.}
\end{quote}

\begin{quote}
The child learned the whole poem by heart.
\end{quote}

Past: \ml{ഹൃദിസ്ഥമാക്കി}, \emph{hṛdisthamākki}.

\ml{ആക്കുക}, to make: to make a thing sit in the heart.

\begin{cousinweb}
\textbf{\ml{ഹൃദിസ്ഥമാക്കുക}} (\emph{hṛdisthamākkuka}) is built from Sanskrit \emph{hṛdistha}, \textquotedblleft{}being in the heart\textquotedblright{} (Monier-Williams), with \ml{ആക്കുക}, \textquotedblleft{}to make\textquotedblright{}, native Dravidian (DEDR 333). English, too, learns by heart.
\end{cousinweb}

\begin{grammarlens}[title={What you've built}]
One more word for giving reasons and opinions.
\end{grammarlens}

\subsection*{Your turn}
\begin{itemize}
\raggedright
  \item \emph{Say it:} \emph{hṛdisthamākkuka}
  \item \emph{Say it:} \emph{hṛdisthamākkuka}, once more
  \item \emph{Say it:} \emph{kuṭṭi ā kavita muḻuvan hṛdisthamākki.} — The child learned the whole poem by heart.
  \item \emph{Recall:} say the Malayalam for studies, then the Malayalam for to grasp, then say \emph{hṛdisthamākkuka} again
\end{itemize}

\begin{tcolorbox}[breakable,colback=teal!4,colframe=teal!35!black,title={Before you move on}]
What does \ml{ഹൃദിസ്ഥമാക്കുക} mean? (\textquotedblleft{}To memorise; to learn by heart\textquotedblright{}.) Say it once more.
\end{tcolorbox}
\section[\texorpdfstring{\ml{പരിശീലിപ്പിക്കുക} (pariśīlippikkuka)}{പരിശീലിപ്പിക്കുക (pariśīlippikkuka)}]{\ml{പരിശീലിപ്പിക്കുക} (pariśīlippikkuka) — to give training; to train}

\label{lesson:ML-C438-parisilippikkuka}



\begin{quote}
Before the new one: say the Malayalam for to grasp, then the Malayalam for to memorise.
\end{quote}

\subsection*{You'll want to know: \ml{പരിശീലിപ്പിക്കുക}}
\textbf{\ml{പരിശീലിപ്പിക്കുക}} — \emph{pariśīlippikkuka} — \textquotedblleft{}to give training; to train\textquotedblright{}.

\begin{quote}
\ml{അവൾ} \ml{പുതിയ} \ml{ജോലിക്കാരെ} \ml{പരിശീലിപ്പിച്ചു}.
\end{quote}

\begin{quote}
\emph{avaḷ putiya jōlikkāre pariśīlippiccu.}
\end{quote}

\begin{quote}
She trained the new staff.
\end{quote}

Past: \ml{പരിശീലിപ്പിച്ചു}, \emph{pariśīlippiccu}.

The causative of \ml{പരിശീലിക്കുക}, to practise: to make someone practise.

\begin{cousinweb}
\textbf{\ml{പരിശീലിപ്പിക്കുക}} (\emph{pariśīlippikkuka}) is built from Sanskrit \emph{pariśīlana}, \textquotedblleft{}contact; application, study\textquotedblright{} (Monier-Williams), with the Malayalam causative verb ending -\ml{ിപ്പിക്കുക}. \ml{പരിശീലനം}, training, is the same word.
\end{cousinweb}

\begin{grammarlens}[title={What you've built}]
One more word for giving reasons and opinions.
\end{grammarlens}

\subsection*{Your turn}
\begin{itemize}
\raggedright
  \item \emph{Say it:} \emph{pariśīlippikkuka}
  \item \emph{Say it:} \emph{pariśīlippikkuka}, once more
  \item \emph{Say it:} \emph{avaḷ putiya jōlikkāre pariśīlippiccu.} — She trained the new staff.
  \item \emph{Recall:} say the Malayalam for to grasp, then the Malayalam for to memorise, then say \emph{pariśīlippikkuka} again
\end{itemize}

\begin{tcolorbox}[breakable,colback=teal!4,colframe=teal!35!black,title={Before you move on}]
What does \ml{പരിശീലിപ്പിക്കുക} mean? (\textquotedblleft{}To give training; to train\textquotedblright{}.) Say it once more.
\end{tcolorbox}
\section[\texorpdfstring{\ml{വ്യാഖ്യാനിക്കുക} (vyākhyānikkuka)}{വ്യാഖ്യാനിക്കുക (vyākhyānikkuka)}]{\ml{വ്യാഖ്യാനിക്കുക} (vyākhyānikkuka) — to interpret}

\label{lesson:ML-C438-vyakhyanikkuka}



\begin{quote}
Before the new one: say the Malayalam for to memorise, then the Malayalam for to give training.
\end{quote}

\subsection*{You'll want to know: \ml{വ്യാഖ്യാനിക്കുക}}
\textbf{\ml{വ്യാഖ്യാനിക്കുക}} — \emph{vyākhyānikkuka} — \textquotedblleft{}to interpret\textquotedblright{}.

\begin{quote}
\ml{ഓരോരുത്തരും} \ml{ആ} \ml{കവിത} \ml{ഓരോ} \ml{വിധത്തിൽ} \ml{വ്യാഖ്യാനിച്ചു}.
\end{quote}

\begin{quote}
\emph{ōrōruttaruṁ ā kavita ōrō vidhattil vyākhyāniccu.}
\end{quote}

\begin{quote}
Everyone interpreted that poem in a different way.
\end{quote}

Past: \ml{വ്യാഖ്യാനിച്ചു}, \emph{vyākhyāniccu}.

The noun is \ml{വ്യാഖ്യാനം}, an interpretation, a commentary.

\begin{cousinweb}
\textbf{\ml{വ്യാഖ്യാനിക്കുക}} (\emph{vyākhyānikkuka}) is built from Sanskrit \emph{vyākhyāna}, \textquotedblleft{}a comment, an explanation\textquotedblright{} (CDIAL 12188), with the Malayalam verb ending -\ml{ിക്കുക}.
\end{cousinweb}

\begin{grammarlens}[title={What you've built}]
One more word for giving reasons and opinions.
\end{grammarlens}

\subsection*{Your turn}
\begin{itemize}
\raggedright
  \item \emph{Say it:} \emph{vyākhyānikkuka}
  \item \emph{Say it:} \emph{vyākhyānikkuka}, once more
  \item \emph{Say it:} \emph{ōrōruttaruṁ ā kavita ōrō vidhattil vyākhyāniccu.} — Everyone interpreted that poem in a different way.
  \item \emph{Recall:} say the Malayalam for to memorise, then the Malayalam for to give training, then say \emph{vyākhyānikkuka} again
\end{itemize}

\begin{tcolorbox}[breakable,colback=teal!4,colframe=teal!35!black,title={Before you move on}]
What does \ml{വ്യാഖ്യാനിക്കുക} mean? (\textquotedblleft{}To interpret\textquotedblright{}.) Say it once more.
\end{tcolorbox}
\section[\texorpdfstring{\ml{നിർവചിക്കുക} (nirvacikkuka)}{നിർവചിക്കുക (nirvacikkuka)}]{\ml{നിർവചിക്കുക} (nirvacikkuka) — to define}

\label{lesson:ML-C438-nirvacikkuka}



\begin{quote}
Before the new one: say the Malayalam for to give training, then the Malayalam for to interpret.
\end{quote}

\subsection*{You'll want to know: \ml{നിർവചിക്കുക}}
\textbf{\ml{നിർവചിക്കുക}} — \emph{nirvacikkuka} — \textquotedblleft{}to define\textquotedblright{}.

\begin{quote}
\ml{അധ്യാപകൻ} \ml{ആ} \ml{വാക്ക്} \ml{ലളിതമായി} \ml{നിർവചിച്ചു}.
\end{quote}

\begin{quote}
\emph{adhyāpakan ā vākkŭ laḷitamāyi nirvaciccu.}
\end{quote}

\begin{quote}
The teacher defined that word in plain words.
\end{quote}

Past: \ml{നിർവചിച്ചു}, \emph{nirvaciccu}.

The noun is \ml{നിർവചനം}, a definition.

\begin{cousinweb}
\textbf{\ml{നിർവചിക്കുക}} (\emph{nirvacikkuka}) is built from Sanskrit \emph{nirvacana}, \textquotedblleft{}speaking out; interpretation, explanation\textquotedblright{} (Monier-Williams), with the Malayalam verb ending -\ml{ിക്കുക}.
\end{cousinweb}

\begin{grammarlens}[title={What you've built}]
One more word for giving reasons and opinions.
\end{grammarlens}

\subsection*{Your turn}
\begin{itemize}
\raggedright
  \item \emph{Say it:} \emph{nirvacikkuka}
  \item \emph{Say it:} \emph{nirvacikkuka}, once more
  \item \emph{Say it:} \emph{adhyāpakan ā vākkŭ laḷitamāyi nirvaciccu.} — The teacher defined that word in plain words.
  \item \emph{Recall:} say the Malayalam for to give training, then the Malayalam for to interpret, then say \emph{nirvacikkuka} again
\end{itemize}

\begin{tcolorbox}[breakable,colback=teal!4,colframe=teal!35!black,title={Before you move on}]
What does \ml{നിർവചിക്കുക} mean? (\textquotedblleft{}To define\textquotedblright{}.) Say it once more.
\end{tcolorbox}
